\documentclass{ximera}

\input{../../../preamble.tex}


\definecolor{fvdnavy}{HTML}{071D43}
\definecolor{fvdcyan}{HTML}{20BFC6}
\definecolor{fvdcyanlight}{HTML}{EFFCFB}
\definecolor{fvdmagenta}{HTML}{F10D69}
\definecolor{fvdmagentalight}{HTML}{FFF1F7}
\definecolor{fvdtext}{HTML}{163044}





%=================================================
% Pedagogische omgevingen
% PDF: tcolorbox
% HTML: learning-box
%=================================================

\newenvironment{voorkennis}
{%
    \ifdefined\HCode
        \HCode{%
<section class="learning-box learning-box--prior">
<div class="learning-box__header">
<span class="learning-box__icon" aria-hidden="true"></span>
<span class="learning-box__title">Voorkennis</span>
</div>
<div class="learning-box__content">
        }%
        \begin{itemize}
    \else
       \begin{tcolorbox}[
    enhanced,
    breakable,
    colback=fvdcyanlight,
    colframe=fvdcyan,
    coltext=fvdtext,
    boxrule=0.5pt,
    leftrule=4pt,
    arc=4mm,
    outer arc=4mm,
    left=7mm,
    right=6mm,
    top=5mm,
    bottom=5mm,
    before skip=6mm,
    after skip=6mm,
    title=Voorkennis,
    coltitle=fvdcyan,
    fonttitle=\bfseries\Large,
    attach title to upper={\par\smallskip},
    boxed title style={
        empty,
        left=0mm,
        right=0mm,
        top=0mm,
        bottom=0mm
    },
    overlay={
        \draw[
            fvdcyan,
            line width=0.5pt
        ]
        ([xshift=42mm,yshift=-13mm]frame.north west)
        --
        ([xshift=-7mm,yshift=-13mm]frame.north east);
    }
]
        \begin{itemize}
    \fi
}
{%
    \end{itemize}
    \ifdefined\HCode
        \HCode{%
</div>
</section>
        }%
    \else
        \end{tcolorbox}
    \fi
}


\newenvironment{leerdoelen}
{%
    \ifdefined\HCode
        \HCode{%
<section class="learning-box learning-box--goals">
<div class="learning-box__header">
<span class="learning-box__icon" aria-hidden="true"></span>
<span class="learning-box__title">Leerdoelen</span>
</div>
<div class="learning-box__content">
        }%
        \begin{itemize}
    \else
\begin{tcolorbox}[
    enhanced,
    breakable,
    colback=fvdmagentalight,
    colframe=fvdmagenta,
    coltext=fvdtext,
    boxrule=0.5pt,
    leftrule=4pt,
    arc=4mm,
    outer arc=4mm,
    left=7mm,
    right=6mm,
    top=5mm,
    bottom=5mm,
    before skip=6mm,
    after skip=6mm,
    title=Leerdoelen,
    coltitle=fvdmagenta,
    fonttitle=\bfseries\Large,
    attach title to upper={\par\smallskip},
    boxed title style={
        empty,
        left=0mm,
        right=0mm,
        top=0mm,
        bottom=0mm
    },
    overlay={
        \draw[
            fvdmagenta,
            line width=0.5pt
        ]
        ([xshift=42mm,yshift=-13mm]frame.north west)
        --
        ([xshift=-7mm,yshift=-13mm]frame.north east);
    }
]
        \begin{itemize}
    \fi
}
{%
    \end{itemize}
    \ifdefined\HCode
        \HCode{%
</div>
</section>
        }%
    \else
        \end{tcolorbox}
    \fi
}


\addPrintStyle{../../..}

\title{Stelsels aan het werk}
\author{Ingmar Herreman}

\begin{abstract}
Een stelsel is meer dan een rij vergelijkingen. Het kan een concrete situatie
beschrijven waarin verschillende voorwaarden tegelijk moeten gelden.

In dit hoofdstuk gebruiken we stelsels als model voor problemen uit onder
andere mengsels, productie en chemie. Het rekenwerk van de matrixreductie
laten we bewust aan ICT over. De nadruk ligt op wat jij moet doen:
onbekenden kiezen, voorwaarden vertalen naar vergelijkingen, een matrix
correct invoeren en de uitkomst interpreteren en controleren.
\end{abstract}

\begin{document}

% FVD_CHAPTER_DATA_START
% Automatisch gegenereerd uit index.tex.
% Niet handmatig aanpassen.
\fvdchapterdata{1}{Van algebra naar ruimte}{2}
% FVD_CHAPTER_DATA_END





\maketitle

\begin{voorkennis}
  \item Je kan een lineair stelsel herkennen en in matrixnotatie schrijven.
  \item Je kan een uitgebreide matrix met ICT reduceren.
  \item Je kan de rang van een matrix uit een gereduceerde vorm aflezen.
  \item Je kan de oplosbaarheid van een lineair stelsel analyseren.
  \item Je kan vrije onbekenden herkennen en een oplossing parametriseren.
\end{voorkennis}

\begin{leerdoelen}
  \item Je kan in een concrete situatie geschikte onbekenden kiezen.
  \item Je kan voorwaarden uit een context vertalen naar lineaire vergelijkingen.
  \item Je kan het verkregen stelsel als uitgebreide matrix invoeren in ICT.
  \item Je kan de ICT-uitvoer interpreteren en terugkoppelen naar de oorspronkelijke context.
  \item Je kan controleren of een wiskundige oplossing ook binnen de context aanvaardbaar is.
  \item Je kan bij een onderbepaald of strijdig model uitleggen wat dit in de context betekent.
  \item Je kan een chemische reactievergelijking balanceren door atoombehoud als een homogeen lineair stelsel te modelleren.
  \item Je kan uitleggen waarom de coëfficiënten van een gebalanceerde reactievergelijking slechts tot op een schaalfactor bepaald zijn.
\end{leerdoelen}

\begin{opmerking}[ICT neemt het rekenwerk over]
In dit hoofdstuk is het niet de bedoeling om omvangrijke matrices volledig
met de hand te reduceren. Gebruik ICT voor de Gauss--Jordan-reductie.

De kern van het werk ligt vóór en na die berekening:
\[
\boxed{\text{context}
\longrightarrow
\text{model}
\longrightarrow
\text{matrix}
\longrightarrow
\text{ICT}
\longrightarrow
\text{interpretatie}}
\]

Een correcte \texttt{rref}-uitvoer is nog geen antwoord op een vraagstuk.
\end{opmerking}


% ============================================================
% 1. VAN CONTEXT NAAR MODEL
% ============================================================

\FVDsection{Van een probleem naar een stelsel}

Bij een gewoon stelsel zijn de onbekenden en vergelijkingen al gegeven.
Bij een toepassing moet je het stelsel eerst \textbf{zelf bouwen}.
Dat noemen we \textbf{mathematiseren}: informatie uit de werkelijkheid
vertalen naar wiskunde.

\FVDsubsection{Een vaste aanpak}

Een betrouwbare aanpak bestaat uit zes stappen.

\begin{enumerate}
  \item \textbf{Lees en structureer.}
        Welke grootheden zijn onbekend? Welke gegevens en voorwaarden zijn er?
  \item \textbf{Kies de onbekenden.}
        Noteer duidelijk wat elke letter betekent, met eenheid waar nodig.
  \item \textbf{Vertaal elke onafhankelijke voorwaarde.}
        Bouw zo een lineair stelsel.
  \item \textbf{Maak de uitgebreide matrix.}
        Bewaak overal dezelfde volgorde van de onbekenden.
  \item \textbf{Reduceer met ICT en analyseer.}
        Lees oplossing, rang en eventuele vrije onbekenden correct af.
  \item \textbf{Keer terug naar de context.}
        Controleer eenheden, teken, gehele aantallen, redelijkheid en de
        oorspronkelijke voorwaarden.
\end{enumerate}

\begin{opmerking}[Een vergelijking moet ergens vandaan komen]
Schrijf niet meteen formules neer. Geef voor jezelf bij elke vergelijking aan
\emph{welke voorwaarde uit de tekst} ze voorstelt. Dat maakt fouten veel
makkelijker zichtbaar.
\end{opmerking}


% ============================================================
% 2. EEN VOLLEDIG VOORBEELD
% ============================================================

\FVDsection{Een volledig modelleerprobleem}

\begin{voorbeeld}[Kippen, eenden en duiven]

Joeri verkoopt zijn kippen, eenden en duiven. Hij heeft in totaal \(27\)
dieren. Het aantal kippen is driemaal het aantal eenden. Een handelaar biedt
\(1\) euro per kip, \(1{,}50\) euro per eend en \(0{,}50\) euro per duif.
Voor alle dieren samen biedt hij \(21\) euro.

Hoeveel dieren heeft Joeri van elke soort?

\textbf{Stap 1 -- onbekenden kiezen}

We nemen
\[
x=\text{aantal kippen},\qquad
y=\text{aantal eenden},\qquad
z=\text{aantal duiven}.
\]

\textbf{Stap 2 -- voorwaarden vertalen}

Uit het totaal aantal dieren:
\[
x+y+z=27.
\]

Uit ``driemaal zoveel kippen als eenden'':
\[
x=3y
\qquad\Longleftrightarrow\qquad
x-3y=0.
\]

Uit het totale bod:
\[
x+1{,}5y+0{,}5z=21.
\]

Dus
\[
\begin{cases}
x+y+z=27,\\
x-3y=0,\\
x+1{,}5y+0{,}5z=21.
\end{cases}
\]

\textbf{Stap 3 -- uitgebreide matrix}

\[
A_b=
\left[
\begin{array}{ccc|c}
1&1&1&27\\
1&-3&0&0\\
1&1{,}5&0{,}5&21
\end{array}
\right].
\]

\textbf{Stap 4 -- reduceren met ICT}

De gereduceerde vorm is
\[
\left[
\begin{array}{ccc|c}
1&0&0&9\\
0&1&0&3\\
0&0&1&15
\end{array}
\right].
\]

Dus
\[
x=9,\qquad y=3,\qquad z=15.
\]

\textbf{Stap 5 -- interpreteren en controleren}

De waarden zijn niet-negatieve gehele aantallen, zoals aantallen dieren
moeten zijn. Bovendien:
\[
9+3+15=27,\qquad 9=3\cdot3,
\]
en
\[
9+1{,}5\cdot3+0{,}5\cdot15=21.
\]

Joeri heeft dus \(9\) kippen, \(3\) eenden en \(15\) duiven.

\end{voorbeeld}


% ============================================================
% 3. ICT
% ============================================================

\FVDsection{Wat doet ICT, en wat doe jij?}

Bij toepassingen is een digitale matrixreductie bijzonder nuttig:
de berekening mag snel gaan, zodat de aandacht naar het model en de
interpretatie kan gaan.

\begin{opmerking}[Jouw verantwoordelijkheid]
ICT kan controleren noch beslissen of
\begin{itemize}
  \item je de juiste onbekenden gekozen hebt;
  \item je de tekst correct naar vergelijkingen vertaald hebt;
  \item de kolommen van de matrix in de juiste volgorde staan;
  \item een negatieve oplossing betekenis heeft;
  \item een kommagetal aanvaardbaar is als je personen of toestellen telt;
  \item een vrije parameter binnen de context alle reële waarden mag aannemen.
\end{itemize}
Dat blijft jouw wiskundige taak.
\end{opmerking}

\FVDsubsection{Een oplossingsverzameling kan door de context kleiner worden}

Stel dat ICT voor een model met drie onbekenden geeft
\[
(x,y,z)=(25+t,\,35-5t,\,4t).
\]

Algebraïsch kan \(t\) elke reële waarde aannemen. Als \(x,y,z\) aantallen
lampen voorstellen, gelden echter extra voorwaarden:
\[
x,y,z\in\mathbb{N}
\qquad\text{en}\qquad
x,y,z\geq0.
\]

De context legt dus beperkingen op die niet automatisch in het oorspronkelijke
lineaire stelsel aanwezig zijn.

\begin{opmerking}
\textbf{Demathematiseren} betekent dat je de wiskundige uitkomst opnieuw
vertaalt naar de werkelijkheid. Een model kan algebraïsch correct opgelost
zijn en toch een onbruikbaar contextantwoord opleveren.
\end{opmerking}


% ============================================================
% 4. MENGPROBLEMEN
% ============================================================

\FVDsection{Mengproblemen}

Mengproblemen leveren natuurlijke lineaire stelsels op omdat verschillende
bestanddelen tegelijk aan opgelegde totalen moeten voldoen.

\begin{voorbeeld}[Voedingsmengsel]

Een bioloog wil een mengsel maken dat \(23\) g proteïne, \(6{,}2\) g vet
en \(16\) g vocht bevat. Hij beschikt over:

\begin{itemize}
  \item mengsel A: \(20\%\) proteïne, \(2\%\) vet en \(15\%\) vocht;
  \item mengsel B: \(10\%\) proteïne, \(6\%\) vet en \(10\%\) vocht;
  \item mengsel C: \(15\%\) proteïne, \(5\%\) vet en \(5\%\) vocht.
\end{itemize}

Neem \(x,y,z\) als het aantal gram van respectievelijk A, B en C.

Dan geeft behoud van elke component:
\[
\begin{cases}
0{,}20x+0{,}10y+0{,}15z=23,\\
0{,}02x+0{,}06y+0{,}05z=6{,}2,\\
0{,}15x+0{,}10y+0{,}05z=16.
\end{cases}
\]

Voer de uitgebreide matrix in ICT in. De reductie levert
\[
x=60,\qquad y=50,\qquad z=40.
\]

De contextcontrole is essentieel. Voor proteïne bijvoorbeeld:
\[
0{,}20\cdot60+0{,}10\cdot50+0{,}15\cdot40
=12+5+6=23.
\]

Er is dus \(60\) g van A, \(50\) g van B en \(40\) g van C nodig.
\end{voorbeeld}


% ============================================================
% 5. RANG IN EEN CONTEXT
% ============================================================

\FVDsection{Wanneer een model niet één antwoord geeft}

De rang vertelt niet alleen iets over een abstract stelsel.
Ook in een context heeft elke mogelijkheid betekenis.

\begin{eigenschap}[Interpretatie van de oplosbaarheid]
Voor een lineair model met \(n\) onbekenden geldt:

\begin{itemize}
  \item \(\operatorname{rang}(A)<\operatorname{rang}(A_b)\):
        de voorwaarden zijn onderling strijdig; er bestaat geen oplossing;
  \item \(\operatorname{rang}(A)=\operatorname{rang}(A_b)=n\):
        de voorwaarden bepalen precies één wiskundige oplossing;
  \item \(\operatorname{rang}(A)=\operatorname{rang}(A_b)<n\):
        de gegeven voorwaarden bepalen de situatie niet volledig;
        er blijven vrije parameters over.
\end{itemize}
\end{eigenschap}

\begin{voorbeeld}[Zestig ledlampen]

Er moeten \(60\) ledlampen worden geplaatst met samen \(40\,000\) lumen.
Er is keuze uit lampen van \(900\), \(500\) en \(400\) lumen.

Met \(x,y,z\) als de aantallen krijgen we
\[
\begin{cases}
x+y+z=60,\\
900x+500y+400z=40\,000.
\end{cases}
\]

Er zijn drie onbekenden maar slechts twee onafhankelijke voorwaarden.
Na reductie blijft één vrije parameter over.

Dat betekent niet dat elke reële parameterwaarde bruikbaar is:
aantallen lampen moeten niet-negatieve gehele getallen zijn.

De context verandert een oneindige algebraïsche oplossingsverzameling dus
in een eindige verzameling praktisch mogelijke keuzes.
\end{voorbeeld}


% ============================================================
% 6. CHEMIE
% ============================================================

\FVDsection{Chemische reactievergelijkingen als stelsel}

Een bijzonder mooie toepassing ontstaat in de chemie.
Bij het balanceren van een reactievergelijking moeten van elk element
evenveel atomen vóór als na de reactie aanwezig zijn.

\begin{opmerking}[Wat wordt behouden?]
Bij het balanceren veranderen we \textbf{alleen de coëfficiënten vóór de
stoffen}. We veranderen nooit de indices in een chemische formule.

De vergelijkingen drukken het behoud van het aantal atomen van elk element
uit. Dat is de wiskundige vertaling die we nodig hebben.
\end{opmerking}

\begin{voorbeeld}[Verbranding van methaan]

We willen
\[
\ldots\,CH_4+\ldots\,O_2
\longrightarrow
\ldots\,CO_2+\ldots\,H_2O
\]
balanceren.

Geef de onbekende coëfficiënten namen:
\[
a\,CH_4+b\,O_2
\longrightarrow
c\,CO_2+d\,H_2O.
\]

Voor elk element schrijven we een behoudsvergelijking.

Koolstof:
\[
a=c.
\]

Waterstof:
\[
4a=2d.
\]

Zuurstof:
\[
2b=2c+d.
\]

Dus:
\[
\begin{cases}
a-c=0,\\
4a-2d=0,\\
2b-2c-d=0.
\end{cases}
\]

De coëfficiënten vormen een homogeen stelsel:
\[
\left[
\begin{array}{cccc|c}
1&0&-1&0&0\\
4&0&0&-2&0\\
0&2&-2&-1&0
\end{array}
\right].
\]

Laat ICT reduceren:
\[
\left[
\begin{array}{cccc|c}
1&0&0&-\frac12&0\\
0&1&0&-1&0\\
0&0&1&-\frac12&0
\end{array}
\right].
\]

Er zijn vier onbekenden en drie pivots:
\[
\operatorname{rang}(A)=3<4.
\]

Er is dus één vrije parameter. Neem \(d=t\). Dan:
\[
a=\frac12t,\qquad b=t,\qquad c=\frac12t,\qquad d=t.
\]

Of:
\[
(a,b,c,d)
=
t\left(\frac12,1,\frac12,1\right).
\]

Chemisch kiezen we de \textbf{kleinste positieve gehele verhouding}.
Met \(t=2\):
\[
(a,b,c,d)=(1,2,1,2).
\]

Daarom:
\[
\boxed{CH_4+2O_2\longrightarrow CO_2+2H_2O}.
\]

\end{voorbeeld}


% ============================================================
% 7. WAAROM EEN VRIJE PARAMETER?
% ============================================================

\FVDsection{Waarom is een reactievergelijking niet uniek?}

Als
\[
CH_4+2O_2\longrightarrow CO_2+2H_2O
\]
klopt, dan klopt ook
\[
2CH_4+4O_2\longrightarrow2CO_2+4H_2O.
\]

Beide beschrijven dezelfde verhouding.

Wiskundig zien we precies hetzelfde:
als \(X\) een oplossing is van het homogene stelsel
\[
AX=0,
\]
dan is voor elke \(\lambda\in\mathbb{R}\)
\[
A(\lambda X)=\lambda AX=0.
\]

Daarom bepaalt het stelsel bij een gewone balanceeropgave geen unieke
coëfficiëntenvector. We zoeken uiteindelijk de kleinste positieve gehele
vertegenwoordiger van die verhouding.

\begin{opmerking}[De nuloplossing]
Omdat het stelsel homogeen is, is
\[
(a,b,c,d)=(0,0,0,0)
\]
altijd een wiskundige oplossing. Chemisch is die betekenisloos.
We zoeken daarom een \textbf{niet-triviale} oplossing met positieve
coëfficiënten.
\end{opmerking}


% ============================================================
% 8. SAMENVATTING
% ============================================================

\FVDsection{Samenvatting: van werkelijkheid naar antwoord}

Bij toepassingen van stelsels is de matrixreductie slechts één schakel.

\[
\boxed{
\begin{array}{c}
\text{1. kies onbekenden}\\
\downarrow\\
\text{2. vertaal voorwaarden naar vergelijkingen}\\
\downarrow\\
\text{3. maak de uitgebreide matrix}\\
\downarrow\\
\text{4. reduceer met ICT}\\
\downarrow\\
\text{5. analyseer rang en oplossingen}\\
\downarrow\\
\text{6. interpreteer en controleer in de context}
\end{array}}
\]

\begin{opmerking}[Wat moet je zeker beheersen?]
Je moet niet alleen een gegeven matrix met ICT kunnen reduceren.
Voor het minimumtraject van dit hoofdstuk moet je ook zelfstandig een
betekenisvolle context kunnen mathematiseren en de verkregen oplossing
correct kunnen demathematiseren.

Bij chemische reactievergelijkingen moet je bovendien het homogene stelsel,
de vrije parameter en de keuze van de kleinste positieve gehele verhouding
kunnen verklaren.
\end{opmerking}

\end{document}
